\section{Conclusion}
\label{sec:conclusion}

This paper has conducted a comprehensive investigation into the role of blockchain technology as an enabler of secure and efficient communication in Unmanned Aerial Vehicle (UAV) networks. We have identified the principal security vulnerabilities and performance bottlenecks inherent in traditional Flying Ad-hoc Networks (FANETs) and critically evaluated the potential of blockchain to address these challenges. Through a detailed review of the state-of-the-art, we have highlighted the significant progress made in areas such as secure routing, decentralized authentication, and lightweight consensus, while also noting the persistent challenges of latency, scalability, and resource overhead.

The central contribution of this work is the proposal and detailed conceptualization of the **PrimeFusion FANET framework**. This novel architecture is distinguished by its synergistic integration of a lightweight, cooperation-based consensus mechanism (Proof-of-Cooperation), an efficient binary data serialization format (CBOR), and a multi-layered security system. Our comparative analysis demonstrates that PrimeFusion offers a more viable path to deploying blockchain in resource-constrained aerial environments compared to the direct application of generic, computationally-intensive blockchain platforms. The framework provides a robust solution for decentralized trust and data integrity while actively minimizing the associated performance penalties.

However, our analysis also maintains a critical perspective, explicitly acknowledging the limitations of the proposed framework, including residual computational overhead, consensus latency, and scalability constraints. These trade-offs underscore that the decision to implement a blockchain-based solution must be carefully weighed against the specific requirements and operational context of the mission. We conclude that while blockchain is not a universal solution, it represents a powerful and transformative technology for a significant class of UAV applications where security, decentralization, and auditable trust are paramount. The future research directions outlined, including the integration of machine learning and the exploration of hierarchical architectures, promise to further enhance the capabilities of blockchain-enabled UAV systems, heralding a new era of secure, autonomous, and cooperative aerial operations.

